\chapter{Nombres complexes}\label{ch:g12:complex}

L'\omterm{def:g10:algebra:equation}{équation} $x^2 = -1$ n'a pas de solution réelle. En élargissant $\R$ d'un
seul nouveau nombre $\iu$ dont le carré vaut $-1$, on obtient le corps $\C$
des \omterm{def:g12:complex:def}{nombres complexes}, dans lequel \emph{toute} \omterm{def:g10:algebra:equation}{équation} polynomiale a des
solutions --- et dont l'arithmétique encode la géométrie du plan~:
translations, rotations et homothéties deviennent additions et
multiplications.

\section{Forme algébrique}

\begin{definition}[Nombres complexes]\label{def:g12:complex:def}
L'ensemble des \emph{nombres complexes}\index{nombre complexe} est
\[
\C = \{a + \iu b : a, b \in \R\},
\]
où $\iu$ est un symbole soumis à la seule règle $\iu^2 = -1$~; l'addition et
la multiplication se font comme pour les polynômes réels en $\iu$, en
réduisant $\iu^2$ à $-1$. Les réels $a$ et $b$ sont la \emph{partie réelle}\index{partie réelle}
$\Rea(z)$ et la \emph{partie imaginaire}\index{partie imaginaire} $\Ima(z)$ de $z = a + \iu b$~;
cette expression est la \emph{forme algébrique}\index{forme algébrique} de $z$, et elle est unique~:
$a + \iu b = a' + \iu b'$ si et seulement si $a = a'$ et $b = b'$.
\end{definition}

\begin{example}
$(2 + 3\iu)(1 - \iu) = 2 - 2\iu + 3\iu - 3\iu^2 = 5 + \iu$.
\end{example}

\begin{definition}[Conjugué]\label{def:g12:complex:conjugate}
Le \emph{conjugué}\index{conjugué} de $z = a + \iu b$ est
$\conj{z} = a - \iu b$.
\end{definition}

\begin{proposition}[Propriétés de la conjugaison]\label{prop:g12:complex:conjugate}
Pour tous $z, w \in \C$~:
\[
\conj{z + w} = \conj z + \conj w, \quad
\conj{zw} = \conj z\,\conj w, \quad
\conj{\conj z} = z, \quad
z + \conj z = 2\Rea(z), \quad
z - \conj z = 2\iu\,\Ima(z),
\]
et $z\conj z = a^2 + b^2 \geq 0$ pour $z = a + \iu b$. De plus $z \in \R$
si et seulement si $z = \conj z$.
\end{proposition}

\begin{proof}
Toutes se vérifient par calcul direct à partir des \omterm{def:g12:complex:def}{formes algébriques}~; par
exemple
$z\conj z = (a + \iu b)(a - \iu b) = a^2 - (\iu b)^2 = a^2 + b^2$.
\end{proof}

\begin{method}[Diviser des nombres complexes]\label{met:g12:complex:division}
Pour écrire un quotient sous \omterm{def:g12:complex:def}{forme algébrique}, multiplier numérateur et
dénominateur par le \omterm{def:g12:complex:conjugate}{conjugué} du dénominateur~:
\[
\frac{1}{2 + \iu} = \frac{2 - \iu}{(2+\iu)(2-\iu)} = \frac{2 - \iu}{5}
= \frac25 - \frac15\iu .
\]
Tout \omterm{def:g12:complex:def}{nombre complexe} non nul a donc un inverse~: $\C$ est un
\emph{corps}.
\end{method}

\section{Le plan complexe, module et argument}

\begin{definition}[Affixe, module]\label{def:g12:complex:modulus}
Dans le plan muni d'un \omterm{def:g10:coordgeom:system}{repère orthonormé}, le point $M(a, b)$ est
l'\emph{\omterm{def:g10:functions:function}{image}} de $z = a + \iu b$, et $z$ est l'\emph{affixe}\index{affixe} de $M$. Le
\emph{module}\index{module} de $z$ est la distance de $M$ à l'origine~:
\[
\abs{z} = \sqrt{a^2 + b^2} = \sqrt{z\conj z}.
\]
\end{definition}

\begin{omfigure}
\begin{tikzpicture}[scale=1.5]
  \draw[->] (-0.7,0) -- (2.9,0) node[right] {$x$};
  \draw[->] (0,-1.9) -- (0,2.1) node[above] {$y$};
  \draw[dashed] (2,1.5) -- (2,0) node[below right] {$a$};
  \draw[dashed] (2,1.5) -- (0,1.5) node[left] {$b$};
  \draw[dashed] (2,-1.5) -- (2,0);
  \draw[-Stealth,omDef,thick] (0,0) -- (2,1.5)
    node[midway, above left, font=\small] {$\abs z$};
  \draw[-Stealth,omThm,thick] (0,0) -- (2,-1.5);
  \fill (2,1.5) circle (1pt) node[above right] {$M(z = a + \iu b)$};
  \fill (2,-1.5) circle (1pt) node[below right] {$\conj z = a - \iu b$};
\end{tikzpicture}

{\small Le point $M(a,b)$ d'\omterm{def:g12:complex:modulus}{affixe} $z$~: le \omterm{def:g12:complex:modulus}{module} $\abs z$ est la
longueur $OM$, et le \omterm{def:g12:complex:conjugate}{conjugué} $\conj z$ est le symétrique de $z$ par
rapport à l'axe réel.}
\end{omfigure}

\begin{proposition}[Propriétés du module]\label{prop:g12:complex:modulus}
Pour $z, w \in \C$~:
\[
\abs{zw} = \abs z\,\abs w, \qquad
\abs{\tfrac zw} = \tfrac{\abs z}{\abs w}\ (w \neq 0), \qquad
\abs{\conj z} = \abs z, \qquad
\abs{z + w} \leq \abs{z} + \abs{w} \ \text{(inégalité triangulaire)}.
\]
De plus $\abs{z_B - z_A}$ est la distance entre les points d'\omterm{def:g12:complex:modulus}{affixes}
$z_A$ et $z_B$.
\end{proposition}

\begin{proof}
$\abs{zw}^2 = zw\,\conj{zw} = z\conj z\, w \conj w = \abs z^2 \abs w^2$
donne la première~; la deuxième et la troisième sont analogues. Pour
l'inégalité triangulaire~:
\[
\abs{z+w}^2 = (z+w)(\conj z + \conj w)
= \abs z^2 + \abs w^2 + 2\Rea(z \conj w)
\leq \abs z^2 + \abs w^2 + 2\abs{z\conj w}
= \bigl(\abs z + \abs w\bigr)^2,
\]
en utilisant $\Rea(u) \leq \abs u$. L'énoncé sur la distance est la formule
de Pythagore appliquée aux \omterm{def:g10:coordgeom:system}{coordonnées} de $z_B - z_A$.
\end{proof}

\begin{definition}[Argument, formes trigonométrique et exponentielle]\label{def:g12:complex:argument}
Soit $z \neq 0$. Un \emph{argument}\index{argument} de $z$, noté
$\Arg(z)$, est toute mesure $\theta$ de l'angle de l'axe réel positif vers
le rayon $OM$~; il est défini à $2k\pi$ près. En notant $r = \abs z$,
\[
z = r(\cos\theta + \iu\sin\theta)
\qquad\text{(forme trigonométrique)}.
\]
En introduisant la notation\index{forme exponentielle}
\[
\eu^{\iu\theta} = \cos\theta + \iu\sin\theta,
\]
cela devient la \emph{forme exponentielle} $z = r\,\eu^{\iu\theta}$.
\end{definition}

\begin{omfigure}
\begin{tikzpicture}[scale=1.5]
  \draw[->] (-0.7,0) -- (2.4,0) node[right] {$x$};
  \draw[->] (0,-0.5) -- (0,1.9) node[above] {$y$};
  \draw[gray, dashed] (1.9,0) arc (0:75:1.9);
  \draw[-Stealth,omDef,thick] (0,0) -- (40:1.9)
    node[pos=0.62, below right=1pt, font=\small] {$r = \abs z$};
  \fill (40:1.9) circle (1pt) node[above right] {$z = r\,\eu^{\iu\theta}$};
  \draw[->,omThm,thick] (0.45,0) arc (0:40:0.45);
  \node[omThm] at (20:0.65) {$\theta$};
\end{tikzpicture}

{\small \omterm{def:g12:complex:argument}{Forme exponentielle}~: $z$ est déterminé par sa distance $r$ à
l'origine et l'angle $\theta$ depuis l'axe réel positif.}
\end{omfigure}

\begin{theorem}\label{thm:g12:complex:expform}
Pour tous $\theta, \theta' \in \R$~:
\[
\eu^{\iu\theta}\,\eu^{\iu\theta'} = \eu^{\iu(\theta + \theta')},
\qquad
\conj{\eu^{\iu\theta}} = \eu^{-\iu\theta},
\qquad
\abs{\eu^{\iu\theta}} = 1 .
\]
Par \omterm{def:g12:seq:sequence}{suite}, pour $z = r\eu^{\iu\theta}$ et $z' = r'\eu^{\iu\theta'}$
non nuls~:
\[
zz' = rr'\,\eu^{\iu(\theta+\theta')},
\qquad
\frac{z}{z'} = \frac{r}{r'}\,\eu^{\iu(\theta - \theta')} :
\]
\emph{les \omterm{def:g12:complex:modulus}{modules} se multiplient, les \omterm{def:g12:complex:argument}{arguments} s'ajoutent.}
\end{theorem}

\begin{proof}
La première identité \emph{est} le couple de formules d'addition
(\cref{prop:g12:trigo:addition})~:
\[
(\cos\theta + \iu\sin\theta)(\cos\theta' + \iu\sin\theta')
= (\cos\theta\cos\theta' - \sin\theta\sin\theta')
+ \iu(\sin\theta\cos\theta' + \cos\theta\sin\theta').
\]
Le reste suit par calcul direct.
\end{proof}

\begin{corollary}[Formule de Moivre]\label{cor:g12:complex:demoivre}
\index{formule de Moivre}
Pour $\theta \in \R$ et $n \in \Z$~:
$\bigl(\cos\theta + \iu\sin\theta\bigr)^n = \cos n\theta + \iu \sin n\theta$.
\end{corollary}

\begin{proof}
Récurrence sur $n \geq 0$ en utilisant \cref{thm:g12:complex:expform}~;
étendre à $n < 0$ en prenant les inverses.
\end{proof}

\begin{method}[Passer d'une forme à l'autre]\label{met:g12:complex:forms}
De l'algébrique à l'exponentielle~: calculer $r = \abs z$, puis trouver
$\theta$ avec $\cos\theta = \frac ar$, $\sin\theta = \frac br$ (situer le
bon quadrant avant d'invoquer un arccosinus). De l'exponentielle à
l'algébrique~: \omterm{def:g10:algebra:expand}{développer} $r\cos\theta + \iu\, r\sin\theta$. Utiliser la
\omterm{def:g12:complex:def}{forme algébrique} pour les sommes, la \omterm{def:g12:complex:argument}{forme exponentielle} pour les produits,
puissances et quotients.
\end{method}

\begin{example}
$z = 1 + \iu$~: $\abs z = \sqrt2$ et $\cos\theta = \sin\theta =
\frac{1}{\sqrt2}$, donc $\theta = \frac\pi4$ et $z = \sqrt2\,\eu^{\iu\pi/4}$.
D'où $z^{20} = 2^{10}\,\eu^{5\iu\pi} = -1024$.
\end{example}

\section{Équations du second degré et racines de l'unité}

\begin{theorem}[Équations du second degré à coefficients réels]\label{thm:g12:complex:quadratic}
Soient $a, b, c \in \R$, $a \neq 0$, et $\Delta = b^2 - 4ac$. L'\omterm{def:g10:algebra:equation}{équation}
$az^2 + bz + c = 0$ a des solutions dans $\C$~:
\[
\begin{aligned}
&\Delta > 0:\ z = \frac{-b \pm \sqrt{\Delta}}{2a} \ (\text{deux réelles}); \qquad
\Delta = 0:\ z = \frac{-b}{2a} \ (\text{double}); \\
&\Delta < 0:\ z = \frac{-b \pm \iu\sqrt{-\Delta}}{2a} \ (\text{deux conjuguées}).
\end{aligned}
\]
\end{theorem}

\begin{proof}
Compléter le carré~:
$az^2 + bz + c = a\left(\left(z + \frac{b}{2a}\right)^2 -
\frac{\Delta}{4a^2}\right)$. Si $\Delta < 0$, alors
$\frac{\Delta}{4a^2} = \left(\iu\frac{\sqrt{-\Delta}}{2a}\right)^2$, et la
différence de carrés se factorise comme d'habitude.
\end{proof}

\begin{definition}[Racines de l'unité]\label{def:g12:complex:rootsofunity}
Pour $n \geq 1$, les \emph{racines $n$-ièmes de l'unité}\index{racines de l'unité}
sont les solutions de $z^n = 1$.
\end{definition}

\begin{theorem}\label{thm:g12:complex:rootsofunity}
L'\omterm{def:g10:algebra:equation}{équation} $z^n = 1$ a exactement $n$ solutions~:
\[
\omega_k = \eu^{2\iu k\pi/n}, \qquad k = 0, 1, \dots, n-1 .
\]
Leurs \omterm{def:g10:functions:function}{images} sont les sommets d'un $n$-gone régulier inscrit dans le cercle
unité, et pour $n \geq 2$ leur somme vaut $0$.
\end{theorem}

\begin{proof}
Chaque $\omega_k$ vérifie $\omega_k^n = \eu^{2\iu k\pi} = 1$, et ils sont
deux à deux distincts car leurs \omterm{def:g12:complex:argument}{arguments} sont dans $\intco{0}{2\pi}$.
Réciproquement, si $z^n = 1$ alors $\abs z^n = 1$, donc $\abs z = 1$ (réel
positif) et $z = \eu^{\iu\theta}$ avec $n\theta \equiv 0 \pmod{2\pi}$,
\emph{c'est-à-dire} $\theta = \frac{2k\pi}{n}$~; en réduisant $k$ modulo $n$
on retombe parmi les $\omega_k$. Les sommets sont espacés de l'angle
$\frac{2\pi}{n}$~: un $n$-gone régulier. Enfin, avec $\omega = \omega_1$,
$\omega \neq 1$ et la somme géométrique donne
\[
\sum_{k=0}^{n-1}\omega_k = \sum_{k=0}^{n-1}\omega^k
= \frac{\omega^n - 1}{\omega - 1} = 0 . \qedhere
\]
\end{proof}

\begin{omfigure}
\begin{tikzpicture}[scale=1.7]
  \draw[->] (-1.3,0) -- (1.45,0) node[right] {$x$};
  \draw[->] (0,-1.3) -- (0,1.4) node[above] {$y$};
  \draw[gray] (0,0) circle (1);
  \draw[omDef] (0:1) -- (72:1) -- (144:1) -- (216:1) -- (288:1) -- cycle;
  \draw[->,omThm,thick] (0.3,0) arc (0:72:0.3);
  \node[omThm, font=\small] at (36:0.52) {$\frac{2\pi}{5}$};
  \fill (0:1) circle (1pt) node[below right] {$1$};
  \fill (72:1) circle (1pt) node[above right] {$\omega$};
  \fill (144:1) circle (1pt) node[above left] {$\omega^2$};
  \fill (216:1) circle (1pt) node[below left] {$\omega^3$};
  \fill (288:1) circle (1pt) node[below right] {$\omega^4$};
\end{tikzpicture}

{\small Les \omterm{def:g11:quad:discriminant}{racines} cinquièmes de l'unité, $\omega = \eu^{2\iu\pi/5}$~: les
sommets d'un pentagone régulier inscrit dans le cercle unité
(\cref{exo:g12:complex:10}).}
\end{omfigure}

\begin{example}
Les racines cubiques de l'unité sont $1$, $j = \eu^{2\iu\pi/3} =
-\frac12 + \iu\frac{\sqrt3}{2}$ et $\conj j = j^2$, vérifiant
$1 + j + j^2 = 0$.
\end{example}

\section{Nombres complexes et géométrie plane}

\begin{proposition}[Dictionnaire géométrique]\label{prop:g12:complex:geometry}
Soient $A, B, C, D$ des points d'\omterm{def:g12:complex:modulus}{affixes} $a, b, c, d$ ($a \neq b$,
$c \neq d$).
\begin{enumerate}
  \item La translation de \omterm{def:g11:vect:vector}{vecteur} d'\omterm{def:g12:complex:modulus}{affixe} $t$ est $z \mapsto z + t$.
  \item La rotation de centre $\omega$ (\omterm{def:g12:complex:modulus}{affixe}) et d'angle $\theta$ est
        $z \mapsto \omega + \eu^{\iu\theta}(z - \omega)$.
  \item L'homothétie de centre $\omega$ et de rapport
        $k \in \R^*$ est $z \mapsto \omega + k(z - \omega)$.
  \item $\dfrac{d - c}{b - a}$ a pour \omterm{def:g12:complex:modulus}{module} $\dfrac{CD}{AB}$ et pour
        \omterm{def:g12:complex:argument}{argument} l'angle entre les \omterm{def:g11:vect:vector}{vecteurs} $\vect{AB}$ et $\vect{CD}$. En
        particulier, $AB \perp CD$ si et seulement si $\frac{d-c}{b-a}$ est
        purement imaginaire, et $A, B, C$ sont alignés ($C \ne A$, $C\ne B$)
        si et seulement si $\frac{c-a}{b-a} \in \R$.
\end{enumerate}
\end{proposition}

\begin{proof}
(1) et (3) reformulent les formules en \omterm{def:g10:coordgeom:system}{coordonnées}. Pour (2), l'application
$w \mapsto \eu^{\iu\theta} w$ multiplie les \omterm{def:g12:complex:modulus}{modules} par $1$ et ajoute
$\theta$ aux \omterm{def:g12:complex:argument}{arguments}~: c'est la rotation d'angle $\theta$ autour de
l'origine~; en conjuguant par la translation envoyant $\omega$ en $0$ on
obtient le centre général. Pour (4), \omterm{def:g12:complex:modulus}{modules} et \omterm{def:g12:complex:argument}{arguments} des quotients se
soustraient (\cref{thm:g12:complex:expform}).
\end{proof}

\section{Exercices}

\begin{exercise}[$\star$]\label{exo:g12:complex:1}
Écrire sous \omterm{def:g12:complex:def}{forme algébrique}~:
$(3 - 2\iu)(1 + 4\iu)$, \quad $\dfrac{2 + \iu}{1 - \iu}$, \quad
$\iu^{2027}$.
\end{exercise}

\begin{exercise}[$\star$]\label{exo:g12:complex:2}
Résoudre dans $\C$~: (a) $z^2 - 4z + 13 = 0$~; (b) $z^2 + z + 1 = 0$.
Vérifier dans chaque cas que les deux solutions sont \omterm{def:g12:complex:conjugate}{conjuguées} et calculer
leur produit.
\end{exercise}

\begin{exercise}[$\star$]\label{exo:g12:complex:3}
Mettre sous \omterm{def:g12:complex:argument}{forme exponentielle}~: $z_1 = -1 + \iu$, $z_2 = \sqrt3 - \iu$,
et calculer $\dfrac{z_1}{z_2}$ sous \omterm{def:g12:complex:argument}{forme exponentielle} puis algébrique. En
déduire la valeur exacte de $\cos\dfrac{11\pi}{12}$.
\end{exercise}

\begin{exercise}[$\star\star$]\label{exo:g12:complex:4}
Décrire géométriquement l'ensemble des points $M$ d'\omterm{def:g12:complex:modulus}{affixe} $z$ tels que~:
(a) $\abs{z - 2} = \abs{z + \iu}$~;\quad
(b) $\abs{z - 1 - \iu} = 3$~;\quad
(c) $\dfrac{z - 1}{z + 1}$ est purement imaginaire.
\end{exercise}

\begin{exercise}[$\star\star$]\label{exo:g12:complex:5}
À l'aide de la formule de Moivre et de la formule du binôme, exprimer
$\cos 3\theta$ comme polynôme en $\cos\theta$, et $\sin 3\theta$ en
\omterm{def:g11:func:function}{fonction} de $\sin\theta$.
\end{exercise}

\begin{exercise}[$\star\star$]\label{exo:g12:complex:6}
\emph{(Linéarisation.)} En utilisant
$\cos\theta = \frac{\eu^{\iu\theta} + \eu^{-\iu\theta}}{2}$, linéariser
$\cos^3\theta$ (l'écrire comme combinaison de $\cos k\theta$), et en déduire
$\displaystyle\int_0^{\pi/2}\cos^3\theta\,\dd\theta$.
\end{exercise}

\begin{exercise}[$\star\star$]\label{exo:g12:complex:7}
Soient $A$ et $B$ d'\omterm{def:g12:complex:modulus}{affixes} $a = 1 + \iu$ et $b = 3 + 2\iu$. Déterminer les
deux points $C$ rendant le triangle $ABC$ équilatéral, comme \omterm{def:g10:functions:function}{images} de $B$
par les rotations de centre $A$ et d'angles $\pm\frac{\pi}{3}$.
\end{exercise}

\begin{exercise}[$\star\star\star$]\label{exo:g12:complex:8}
Résoudre $z^4 = -4$ dans $\C$ et placer les solutions. Factoriser
$z^4 + 4$ en deux polynômes du second degré à coefficients réels.
\end{exercise}

\begin{exercise}[$\star\star\star$]\label{exo:g12:complex:9}
Pour $\theta \in \intoo{0}{2\pi}$ et $n \in \N$, calculer
\[
S = 1 + \cos\theta + \cos 2\theta + \dots + \cos n\theta
\]
en sommant la progression géométrique
$\sum_{k=0}^n \eu^{\iu k\theta}$ et en prenant les \omterm{def:g12:complex:def}{parties réelles}.
(Réponse~:
$S = \dfrac{\sin\frac{(n+1)\theta}{2}}{\sin\frac{\theta}{2}}
\cos\frac{n\theta}{2}$.)
\end{exercise}

\begin{exercise}[$\star\star\star$]\label{exo:g12:complex:10}
Soit $\omega = \eu^{2\iu\pi/5}$.
\begin{enumerate}
  \item Justifier que $1 + \omega + \omega^2 + \omega^3 + \omega^4 = 0$.
  \item Soient $u = \omega + \omega^4$ et $v = \omega^2 + \omega^3$. Montrer
        que $u + v = -1$ et $uv = -1$, et en déduire que
        $u = \frac{-1+\sqrt5}{2}$.
  \item Conclure que $\cos\dfrac{2\pi}{5} = \dfrac{\sqrt5 - 1}{4}$.
\end{enumerate}
\end{exercise}

\section{Problème~: le pentagone dans le miroir de l'unité}

\begin{problem}[{Devoir du week-end --- les racines cinquièmes de
l'unité calculent $\cos 72^\circ$ exactement, le nombre d'or signe le
résultat, et la multiplication se révèle être une rotation}]
\label{pb:g12:complex:1}
Euclide savait construire le pentagone régulier, mais la raison
\emph{pour laquelle} il cède à la règle et au compas --- alors que
l'humble angle de $20^\circ$ y résiste --- est restée cachée deux
mille ans. Elle se cache dans ce chapitre~: les cinq \omterm{def:g11:quad:discriminant}{racines}
cinquièmes de l'unité (\cref{thm:g12:complex:rootsofunity}) ont une
somme nulle, et cette seule ligne d'algèbre extrait $\cos 72^\circ$
d'une \omterm{def:g10:algebra:equation}{équation} du second degré --- des racines carrées seulement,
donc constructible, et signé, bien entendu, par le nombre d'or.
Autour de ce joyau~: de l'aisance, le raccourci de Moivre vers la
trigonométrie, et la multiplication démasquée en géométrie.

\medskip
\noindent\textbf{Partie I --- Aisance.}
\begin{enumerate}
  \item Calculer $(2 + \iu)(3 - \iu)$, \
        $\dfrac{1 + \iu}{1 - \iu}$
        (\cref{met:g12:complex:division}) et $\abs{3 + 4\iu}$.
  \item Mettre sous \omterm{def:g12:complex:argument}{forme exponentielle}~: $1 + \iu$~; \ $-2$~; \
        $1 - \iu\sqrt3$.
  \item Calculer $(1 + \iu)^8$ à l'aide de la \omterm{def:g12:complex:argument}{forme
        exponentielle}.
  \item Résoudre $z^2 = \iu$, puis $z^2 - 2z + 5 = 0$
        (\cref{thm:g12:complex:quadratic}).
  \item Dictionnaire géométrique
        (\cref{prop:g12:complex:geometry})~: décrire
        l'application $z \mapsto \iu z$, ainsi que l'ensemble
        $\abs{z - (1 + \iu)} = 2$.
\end{enumerate}

\medskip
\noindent\textbf{Partie II --- Cinq \omterm{def:g11:quad:discriminant}{racines}, un pentagone.}
Posons $\omega = \eu^{2\iu\pi/5}$ et considérons
$1, \omega, \omega^2, \omega^3, \omega^4$.
\begin{enumerate}[resume]
  \item Vérifier que ce sont exactement les solutions de
        $z^5 = 1$, et décrire la figure qu'elles dessinent dans le
        plan.
  \item Démontrer que leur somme est nulle. (\omterm{def:g11:seq:geometric}{Suite géométrique} ---
        ou bien multiplier la somme par $1 - \omega$.)
  \item Prendre les \omterm{def:g12:complex:def}{parties réelles} et en déduire
        \[
        1 + 2\cos\frac{2\pi}{5} + 2\cos\frac{4\pi}{5} = 0 .
        \]
  \item Poser $c = \cos\frac{2\pi}{5}$ et utiliser la formule de
        duplication pour $\cos\frac{4\pi}{5}$ afin de convertir la
        question 8 en une \omterm{def:g10:algebra:equation}{équation} du second degré en $c$~; la
        résoudre et conclure
        \[
        \cos 72^\circ = \frac{\sqrt5 - 1}{4} .
        \]
  \item Vérifier numériquement, puis laisser le nombre d'or signer
        son œuvre~: montrer que $c = \frac{1}{2\varphi}$, où
        $\varphi = \frac{1 + \sqrt5}{2}$
        (\cref{pb:g10:algebra:1}) --- et énoncer le célèbre écho
        géométrique~: dans un pentagone régulier, la diagonale vaut
        $\varphi$ fois le côté.
\end{enumerate}

\medskip
\noindent\textbf{Partie III --- Les raccourcis de Moivre.}
\begin{enumerate}[resume]
  \item Le verdict de constructibilité~: $\cos 72^\circ$ ne réclame
        que des racines carrées (question 9), donc le pentagone est
        constructible à la règle et au compas~; $\cos 20^\circ$
        vérifie une \omterm{def:g10:algebra:equation}{équation} du troisième degré irréductible
        (\cref{pb:g12:trigo:1}), donc l'angle de $60^\circ$ ne peut
        être trisecté. Énoncer le critère qui se dessine (les
        racines carrées se construisent, les racines cubiques non),
        dont la théorie complète --- Gauss, à dix-neuf ans, et le
        polygone à 17 côtés --- réside dans les volumes
        universitaires.
  \item Résoudre $z^4 = -4$ par la \omterm{def:g12:complex:argument}{forme exponentielle} et retrouver
        la factorisation de l'\cref{exo:g12:complex:8},
        $z^4 + 4 = (z^2 - 2z + 2)(z^2 + 2z + 2)$, en appariant les
        \omterm{def:g11:quad:discriminant}{racines} \omterm{def:g12:complex:conjugate}{conjuguées}.
  \item Développer $(\cos\theta + \iu\sin\theta)^3$ par la formule
        du binôme et par la formule de Moivre, et retrouver
        $\cos 3\theta = 4\cos^3\theta - 3\cos\theta$ en trois lignes
        --- comparer avec la voie des formules d'addition du
        \cref{pb:g12:trigo:1}.
  \item Recouper avec l'\cref{exo:g12:complex:9}~: pour
        $\theta = \frac{2\pi}{5}$ et $n = 4$, que vaut la somme
        $1 + \cos\theta + \dots + \cos 4\theta$, et pourquoi
        est-ce cohérent avec la question 8~?
  \item Calculer les trois racines cubiques de $8\iu$ sous \omterm{def:g12:complex:def}{forme
        algébrique}.
\end{enumerate}

\medskip
\noindent\textbf{Partie IV --- La multiplication est une
géométrie.}
\begin{enumerate}[resume]
  \item Décrire complètement l'application
        $z \mapsto (1 + \iu)z$~: de quel angle fait-elle tourner,
        de quel facteur dilate-t-elle, et que devient le carré
        unité~?
  \item Démontrer $\abs{z_1 z_2} = \abs{z_1}\,\abs{z_2}$ à l'aide
        des \omterm{def:g12:complex:conjugate}{conjugués} (\cref{prop:g12:complex:modulus}), puis
        calculer $(1+\iu)^n$ pour $n = 1, 2, 3, 4$ et décrire la
        spirale que tracent les puissances.
  \item Le nombre $\mathrm{j} = \eu^{2\iu\pi/3}$~: montrer que
        $1 + \mathrm{j} + \mathrm{j}^2 = 0$, et vérifier sur le
        triangle $\left(1, \mathrm{j}, \mathrm{j}^2\right)$ le
        critère classique~: $a + \mathrm{j} b + \mathrm{j}^2 c = 0$
        caractérise (pour une orientation) un triangle équilatéral
        $abc$.
  \item Le théorème fondamental de l'algèbre (admis~:
        d'Alembert--Gauss)~: tout polynôme se factorise complètement
        sur $\C$. En déduire le corollaire pour le monde réel ---
        tout polynôme réel se factorise en facteurs réels du premier
        et du second degré --- et citer la question 12 comme
        exemple.
  \item Pour finir --- le portrait de $\C$, une phrase pour
        chacune~: les nombres deviennent des points~; la
        multiplication devient rotation et dilatation~; les \omterm{def:g12:complex:rootsofunity}{racines
        de l'unité} deviennent des polygones réguliers~; Moivre
        transforme les puissances en identités trigonométriques~; et
        la clôture algébrique garantit qu'aucune \omterm{def:g10:algebra:equation}{équation} n'aura
        jamais besoin d'un monde plus vaste. Coda~: quels deux
        nombres célèbres se sont rencontrés à la question 10~?
\end{enumerate}
\end{problem}
